% TOP_PROBLEM: {"id": 2511, "title": "Kourovka Notebook Problem 21.2", "category_id": 17, "category": {"id": 17, "name": "group_theory", "display_name": "Group Theory", "description": "Problems about groups, group actions, representations, and related algebraic structures.", "slug": "group-theory", "order_index": 17, "created_at": "2026-07-31T15:26:25.671Z"}, "status": "open", "problem_number": "KOU-21.2", "set_id": 10}
% FIRST_POSTED: 2026-10-01
% ENTRY_KIND: statement_only
% UnsolvedMath ID 2511; source catalogue code KOU-21.2
% !TeX program = lualatex
% Definition and sources only; this file is not an LLM solution attempt.
\documentclass[11pt,a4paper]{article}
\usepackage[margin=25mm]{geometry}
\usepackage{fontspec}
\setmainfont{lmroman10-regular.otf}[BoldFont=lmroman10-bold.otf,
 ItalicFont=lmroman10-italic.otf,BoldItalicFont=lmroman10-bolditalic.otf]
\usepackage{amsmath,amssymb,mathtools}
\usepackage{unicode-math}
\setmathfont{latinmodern-math.otf}
\AtBeginDocument{\let\setminus\smallsetminus}
\usepackage{xurl,hyperref}
\hypersetup{colorlinks=true,urlcolor=blue}
\setcounter{secnumdepth}{0}
\setlength{\parindent}{0pt}
\setlength{\parskip}{5pt}
\setlength{\emergencystretch}{3em}
\begin{document}
\section*{Kourovka Notebook Problem 21.2}\label{2511}
{\small UnsolvedMath ID 2511 · KOU-21.2 · Group Theory · Kourovka Notebook - New Problems, Issue 21}
\subsection{Definitions and mathematical statement}
For an element $x$ of a finite group, let $o(x)$ be its order, the least
positive integer $m$ with $x^m=1$. A simple group is a nontrivial group
whose only normal subgroups are $1$ and the whole group; normality means
invariance under conjugation by every group element. Cyclic groups of
prime order are included in this definition.

Let $S$ be finite simple and $G$ any finite group. Suppose there exists
a bijection of their underlying sets $f:G\to S$ such that
\[
 o(x)\mid o(f(x))\qquad\text{for every }x\in G.
\]
Here divisibility is divisibility of positive integers. The question is
whether $G$ must itself be simple.

The map $f$ is not required to preserve multiplication, inverses,
conjugacy classes, or the identity. In fact the hypotheses imply
$f(1_G)=1_S$: the preimage of $1_S$ must have order dividing one.
The bijection also implies $|G|=|S|$, but the requested conclusion is
simplicity, not an asserted isomorphism $G\cong S$.

The condition matches individual elements without repetitions: each
element of $S$ can serve as the image of only one element of $G$.
It is stronger than merely demanding that every element order in $G$
divide some element order in $S$, because multiplicities must be matched.
A negative answer would require a nonsimple group of the same order as
some simple group together with such a complete matching of elements.
\subsection{Short English statement}
If the elements of a finite group can be bijectively matched to those of a simple group with each order dividing its matched order, must the first group be simple?
\subsection{Sources}
\begin{itemize}
\item[S1] ULAM.AI and the UnsolvedMath Contributors, supplied problems.json, record 2511 (KOU-21.2), Kourovka Notebook - New Problems, Issue 21. Catalogue identity and source transcription; CC BY 4.0.
\url{https://huggingface.co/datasets/ulamai/UnsolvedMath}.
\item[S2] The Kourovka Notebook, 21st edition, new problems, Problem 21.2. Expanded formulation of the supplied catalogue statement.
\url{https://arxiv.org/pdf/1401.0300v47}.
\end{itemize}
\end{document}
